\documentclass[12pt]{article}
\usepackage{amsmath, amssymb, hyperref, geometry, booktabs}
\geometry{margin=2.5cm}

\title{Resúmenes Estructurados: Máquinas de Boltzmann y Mecánica Estadística}
\author{Preparado para el curso de Mecánica Estadística --- IFA, Universidad de Valparaíso}
\date{\today}

\begin{document}
\maketitle

\noindent Este documento resume 11 papers sobre Máquinas de Boltzmann, modelos basados en
energía y mecánica estadística, organizados en torno a cuatro ejes temáticos: (1) la
Distribución de Boltzmann, (2) pesos probabilísticos y modelos basados en energía, (3) la
función de partición $Z$, y (4) ensembles como lenguaje de inferencia.

\bigskip
\hrule
\bigskip

%=================================================================
\section{Jaynes (1957a) --- Information Theory and Statistical Mechanics}

\subsection*{Referencia completa}
E. T. Jaynes (1957). \textit{Information Theory and Statistical Mechanics}. Physical Review,
106(4), 620--630. DOI: \href{https://doi.org/10.1103/PhysRev.106.620}{10.1103/PhysRev.106.620}.
\cite{jaynes1957a}

\subsection*{Pregunta central}
¿Puede la mecánica estadística de equilibrio derivarse como un caso particular de inferencia
estadística (máxima entropía), sin presuponer hipótesis físicas adicionales como la
ergodicidad o la equiprobabilidad a priori?

\subsection*{Aporte principal}
Jaynes demuestra que maximizar la entropía de Shannon sujeta a los valores esperados conocidos
(p.\ ej.\ la energía) reproduce exactamente las reglas de cálculo de la mecánica estadística de
Gibbs, incluida la función de partición. Reinterpreta así la mecánica estadística como un caso
de inferencia bayesiana --- ``mecánica estadística subjetiva'' --- eliminando la necesidad de
supuestos ergódicos.

\subsection*{Conexión con los ejes temáticos}
\begin{itemize}
  \item \textbf{Distribución de Boltzmann:} Sí --- es la derivación original vía máxima entropía; artículo fundacional del eje.
  \item \textbf{Pesos probabilísticos y EBM:} No aplica directamente --- no trata redes neuronales, pero establece el marco $P(s) \propto e^{-\beta E(s)}$ que el aprendizaje automático reinterpretará décadas después como EBM.
  \item \textbf{Función de partición $Z$:} Sí --- introduce $Z(\lambda)$ explícitamente como normalizador derivado de los multiplicadores de Lagrange.
  \item \textbf{Ensembles como inferencia:} Sí --- es el artículo fundacional de esta equivalencia.
\end{itemize}

\subsection*{Ecuaciones o resultados clave}
\begin{equation}
p_i = \exp\{-[\lambda_0+\lambda_1 f_1(x_i)+\cdots+\lambda_m f_m(x_i)]\}
\end{equation}
\begin{equation}
Z(\lambda_1,\ldots,\lambda_m) = \sum_i \exp\{-[\lambda_1 f_1(x_i)+\cdots+\lambda_m f_m(x_i)]\}, \qquad
\langle f(x)\rangle = -\frac{\partial \ln Z}{\partial \lambda}
\end{equation}

\subsection*{Cita textual}
\begin{quote}
\textit{``we must use that probability distribution which has maximum entropy''}
\end{quote}

\subsection*{Relevancia para una presentación de 10 minutos}
\begin{tabular}{@{}p{2cm}p{11.5cm}@{}}
\toprule
\textbf{Nivel} & \textbf{Justificación} \\
\midrule
Alta & Es el origen conceptual de los cuatro ejes temáticos del proyecto; debería abrir la presentación como marco unificador. \\
\bottomrule
\end{tabular}

\bigskip
\hrule
\bigskip

%=================================================================
\section{Jaynes (1957b) --- Information Theory and Statistical Mechanics II}

\subsection*{Referencia completa}
E. T. Jaynes (1957). \textit{Information Theory and Statistical Mechanics II}. Physical
Review, 108(2), 171--190. DOI: \href{https://doi.org/10.1103/PhysRev.108.171}{10.1103/PhysRev.108.171}.
\cite{jaynes1957b}

\subsection*{Pregunta central}
¿Cómo se extiende el principio de máxima entropía al formalismo de matriz de densidad
(sistemas cuánticos), y qué relación guarda con la irreversibilidad y la pérdida de
información?

\subsection*{Aporte principal}
Extiende el tratamiento de inferencia estadística de la Parte I al formalismo cuántico de la
matriz de densidad, mostrando que la matriz de densidad de máxima entropía sujeta a
restricciones adopta la forma de Boltzmann-Gibbs cuántica. Introduce el principio de
``complementariedad estadística'': las probabilidades verificables empíricamente corresponden
necesariamente a predicciones incompletas, vinculando irreversibilidad con pérdida de
información.

\subsection*{Conexión con los ejes temáticos}
\begin{itemize}
  \item \textbf{Distribución de Boltzmann:} Sí --- versión cuántica vía matriz de densidad, $\rho \propto e^{-\beta H}$.
  \item \textbf{Pesos probabilísticos y EBM:} No aplica directamente --- es física cuántica formal, no machine learning, aunque la entropía de von Neumann reaparece en EBMs variacionales modernos.
  \item \textbf{Función de partición $Z$:} Sí --- $Z = \mathrm{Tr}(e^{-\beta H})$ en el formalismo de operadores.
  \item \textbf{Ensembles como inferencia:} Sí --- generaliza el argumento de inferencia a sistemas no estacionarios e irreversibles.
\end{itemize}

\subsection*{Ecuaciones o resultados clave}
\begin{equation}
S = -\mathrm{Tr}(\rho \ln \rho)
\end{equation}
\begin{equation}
\langle F \rangle = \sum_i w_i \langle i|F|i\rangle = \mathrm{Tr}(\rho F)
\end{equation}

\subsection*{Cita textual}
\begin{quote}
\textit{``the predictive aspect of statistical mechanics as a form of statistical inference''}
\end{quote}

\subsection*{Relevancia para una presentación de 10 minutos}
\begin{tabular}{@{}p{2cm}p{11.5cm}@{}}
\toprule
\textbf{Nivel} & \textbf{Justificación} \\
\midrule
Media & Profundiza el marco teórico, pero el formalismo cuántico de matriz de densidad excede el alcance de una charla introductoria; útil como referencia de rigor, no como contenido central. \\
\bottomrule
\end{tabular}

\bigskip
\hrule
\bigskip

%=================================================================
\section{Ackley, Hinton \& Sejnowski (1985) --- A Learning Algorithm for Boltzmann Machines}

\subsection*{Referencia completa}
D. H. Ackley, G. E. Hinton \& T. J. Sejnowski (1985). \textit{A Learning Algorithm for
Boltzmann Machines}. Cognitive Science, 9(1), 147--169. DOI:
\href{https://doi.org/10.1207/s15516709cog0901_7}{10.1207/s15516709cog0901\_7}. \cite{ackley1985}

\subsection*{Pregunta central}
¿Puede una red estocástica de unidades binarias con conexiones simétricas aprender, mediante
una regla de actualización completamente local, a reproducir una distribución de probabilidad
objetivo sobre sus unidades visibles?

\subsection*{Aporte principal}
Introduce la Máquina de Boltzmann: una red cuya dinámica estocástica converge al equilibrio
térmico (distribución de Boltzmann). Deriva una regla de aprendizaje basada en el gradiente de
una medida de divergencia $G$ (equivalente a la divergencia de Kullback-Leibler) que depende
solo de correlaciones locales $\langle s_i s_j\rangle$, calculables comparando una fase
``despierta'' (datos fijos) y una fase ``dormida'' (equilibrio libre).

\subsection*{Conexión con los ejes temáticos}
\begin{itemize}
  \item \textbf{Distribución de Boltzmann:} Sí --- es la base física de la dinámica de la red: $P_\alpha \propto e^{-E_\alpha/T}$.
  \item \textbf{Pesos probabilísticos y EBM:} Sí --- define explícitamente la energía y la interpretación de $w_{ij}$ como acoplamientos aprendibles; es el eje central del paper.
  \item \textbf{Función de partición $Z$:} Implícita --- $Z$ normaliza $P_\alpha$, pero el algoritmo está diseñado precisamente para \emph{evitar} calcularla explícitamente, usando diferencias de correlaciones en su lugar.
  \item \textbf{Ensembles como inferencia:} Sí --- el aprendizaje se plantea como minimización de una divergencia entre la distribución de los datos y la del modelo, precursor directo de la lectura bayesiana moderna.
\end{itemize}

\subsection*{Ecuaciones o resultados clave}
\begin{equation}
E = -\sum_{i<j} w_{ij}\,s_i s_j + \sum_i \theta_i\, s_i
\end{equation}
\begin{equation}
\Delta w_{ij} = \varepsilon\left(\langle s_i s_j\rangle_{\text{datos}} - \langle s_i s_j\rangle_{\text{libre}}\right)
\end{equation}

\subsection*{Cita textual}
\begin{quote}
\textit{``a general learning rule for modifying the connection strengths''}
\end{quote}

\subsection*{Relevancia para una presentación de 10 minutos}
\begin{tabular}{@{}p{2cm}p{11.5cm}@{}}
\toprule
\textbf{Nivel} & \textbf{Justificación} \\
\midrule
Alta & Artículo fundacional de la Máquina de Boltzmann; conecta directamente energía, pesos y aprendizaje. Imprescindible. \\
\bottomrule
\end{tabular}

\bigskip
\hrule
\bigskip

%=================================================================
\section{LeCun, Chopra, Hadsell, Ranzato \& Huang (2006) --- A Tutorial on Energy-Based Learning}

\subsection*{Referencia completa}
Y. LeCun, S. Chopra, R. Hadsell, M. Ranzato \& F. J. Huang (2006). \textit{A Tutorial on
Energy-Based Learning}. En: \textit{Predicting Structured Data}, MIT Press.
\href{http://yann.lecun.com/exdb/publis/pdf/lecun-06.pdf}{PDF}. \cite{lecun2006}

\subsection*{Pregunta central}
¿Cómo puede formularse el aprendizaje automático --- probabilístico y no probabilístico ---
dentro de un único marco unificado basado en funciones de energía, evitando el problema de la
normalización?

\subsection*{Aporte principal}
Generaliza la distribución de Boltzmann a una clase amplia de modelos de aprendizaje (EBMs),
donde la inferencia es minimización de energía y el aprendizaje consiste en asignar energías
bajas a configuraciones observadas. Muestra que los modelos probabilísticos son un caso
particular de EBM (cuando se exige normalización), liberando el diseño de arquitecturas de la
necesidad de calcular $Z$.

\subsection*{Conexión con los ejes temáticos}
\begin{itemize}
  \item \textbf{Distribución de Boltzmann:} Sí --- $p(x) \propto e^{-E(x)}$ es el punto de partida explícito de los EBM.
  \item \textbf{Pesos probabilísticos y EBM:} Sí --- es la referencia central que define formalmente los EBM y la energía parametrizada $E_\theta$.
  \item \textbf{Función de partición $Z$:} Sí --- discute extensamente la intratabilidad de $Z$ y cómo el enfoque EBM la evita por diseño.
  \item \textbf{Ensembles como inferencia:} Parcial --- discute la función de pérdida como sustituto de $-\ln$ verosimilitud, pero no enmarca explícitamente el aprendizaje como inferencia bayesiana/ensembles.
\end{itemize}

\subsection*{Ecuaciones o resultados clave}
\begin{equation}
P(Y\mid X) = \frac{e^{-\beta E(Y,X)}}{\int_{y} e^{-\beta E(y,X)}\,dy}
\end{equation}
\begin{equation}
W \leftarrow W - \eta\left(\frac{\partial E(Y_i,X_i,W)}{\partial W} - \frac{\partial E(\bar Y_i,X_i,W)}{\partial W}\right)
\end{equation}

\subsection*{Cita textual}
\begin{quote}
\textit{``associating a scalar energy to each configuration of the variables''}
\end{quote}

\subsection*{Relevancia para una presentación de 10 minutos}
\begin{tabular}{@{}p{2cm}p{11.5cm}@{}}
\toprule
\textbf{Nivel} & \textbf{Justificación} \\
\midrule
Alta & Formaliza por completo el eje de modelos basados en energía; es la pieza central para conectar física e IA moderna. \\
\bottomrule
\end{tabular}

\bigskip
\hrule
\bigskip

%=================================================================
\section{Cheng, Chen \& Wang (2018) --- Boltzmann Machines versus Born Machines}

\subsection*{Referencia completa}
S. Cheng, J. Chen \& L. Wang (2018). \textit{Information Perspective to Probabilistic
Modeling: Boltzmann Machines versus Born Machines}. Entropy, 20(8), 583. DOI:
\href{https://doi.org/10.3390/e20080583}{10.3390/e20080583}. \cite{cheng2018}

\subsection*{Pregunta central}
¿Qué diferencias en capacidad expresiva existen entre los modelos generativos basados en
energía (Máquinas de Boltzmann) y los basados en estados cuánticos (Máquinas de Born), y cómo
lo revela la teoría de la información clásica frente a la cuántica?

\subsection*{Aporte principal}
Compara las RBM ($p(v)=e^{-E(v)}/Z$) con las ``Born machines'' ($p(v)=|\Psi(v)|^2/N$) usando
cotas de información mutua y entropía de entrelazamiento de Rényi. Muestra que diseñar
arquitecturas (conexiones locales/dispersas) que respeten el patrón de información mutua del
dataset mejora la eficiencia de aprendizaje en ambos enfoques, validado en MNIST.

\subsection*{Conexión con los ejes temáticos}
\begin{itemize}
  \item \textbf{Distribución de Boltzmann:} Sí --- el RBM clásico es el objeto de comparación directa.
  \item \textbf{Pesos probabilísticos y EBM:} Sí --- analiza explícitamente la arquitectura de pesos $w_{ij}$ del RBM y su relación con la información mutua de los datos.
  \item \textbf{Función de partición $Z$:} Sí --- $Z$ aparece en la normalización del RBM y se discute su intratabilidad como motivación de los enfoques cuánticos alternativos.
  \item \textbf{Ensembles como inferencia:} Indirecta --- no usa el lenguaje de ensembles/Bayes; se centra en teoría de la información clásica y cuántica.
\end{itemize}

\subsection*{Ecuaciones o resultados clave}
\begin{equation}
p(v) = \frac{e^{-E(v)}}{Z}, \qquad Z = \sum_v e^{-E(v)}
\end{equation}
\begin{equation}
p(v) = \frac{|\Psi(v)|^2}{N} \qquad \text{(Máquina de Born)}
\end{equation}

\subsection*{Cita textual}
\begin{quote}
\textit{``have the potential to be good generative models for machine learning''}
\end{quote}

\subsection*{Relevancia para una presentación de 10 minutos}
\begin{tabular}{@{}p{2cm}p{11.5cm}@{}}
\toprule
\textbf{Nivel} & \textbf{Justificación} \\
\midrule
Media & Aporta una conexión interesante entre los ejes 1 y 3, pero el contenido cuántico (entrelazamiento, Born machines) es denso para 10 minutos; mejor como nota comparativa breve. \\
\bottomrule
\end{tabular}

\bigskip
\hrule
\bigskip

%=================================================================
\section{Pu\v{s}karov \& Cort\'es Cubero (2018) --- Machine Learning Algorithms Based on Generalized Gibbs Ensembles}

\subsection*{Referencia completa}
T. Pu\v{s}karov \& A. Cort\'es Cubero (2018). \textit{Machine Learning Algorithms Based on
Generalized Gibbs Ensembles}. arXiv:1804.03546 [cond-mat.stat-mech]. \cite{puskarov2018}

\subsection*{Pregunta central}
¿Puede sustituirse la máquina de Boltzmann cuántica --- de entrenamiento ineficiente debido a
la no conmutatividad de operadores --- por un algoritmo basado en el ensemble de Gibbs
generalizado (GGE), que sí pueda entrenarse eficientemente por descenso de gradiente?

\subsection*{Aporte principal}
Propone un algoritmo de aprendizaje basado en el GGE, que reemplaza la única temperatura
inversa $\beta$ por un conjunto de ``temperaturas efectivas'' $\beta_n$ asociadas a cargas
conservadas de un Hamiltoniano integrable, y demuestra que es la única máquina de Boltzmann
cuántica entrenable eficientemente, gracias a que las cargas conservadas conmutan entre sí. Su
versión clásica simplificada clasifica MNIST con $\sim$36 veces menos parámetros que una RBM
comparable.

\subsection*{Conexión con los ejes temáticos}
\begin{itemize}
  \item \textbf{Distribución de Boltzmann:} Sí --- el GGE es una generalización directa del ensemble canónico/distribución de Boltzmann.
  \item \textbf{Pesos probabilísticos y EBM:} Sí --- las ``temperaturas efectivas'' $\beta_n$ desempeñan un papel análogo a los pesos $w_{ij}$.
  \item \textbf{Función de partición $Z$:} Sí --- $Z=\mathrm{Tr}(e^{-\sum_n \beta_n Q_n})$ generaliza la función de partición canónica a múltiples cargas conservadas.
  \item \textbf{Ensembles como inferencia:} Sí --- es el ejemplo más directo de extender el diccionario ensemble$\leftrightarrow$ML del eje 4 más allá del ensemble canónico.
\end{itemize}

\subsection*{Ecuaciones o resultados clave}
\begin{equation}
P_\psi = \frac{e^{-\sum_n \beta_n Q_n}}{Z}, \qquad Z = \mathrm{Tr}\!\left(e^{-\sum_n \beta_n Q_n}\right)
\end{equation}
\begin{equation}
\mathcal{L} = S(\rho_{\text{data}}\|\rho) = \mathrm{Tr}[\rho_{\text{data}}\log\rho_{\text{data}}] - \mathrm{Tr}[\rho_{\text{data}}\log\rho]
\end{equation}

\subsection*{Cita textual}
\begin{quote}
\textit{``GGE algorithm turns out to be an optimal implementation of a quantum Boltzmann machine''}
\end{quote}

\subsection*{Relevancia para una presentación de 10 minutos}
\begin{tabular}{@{}p{2cm}p{11.5cm}@{}}
\toprule
\textbf{Nivel} & \textbf{Justificación} \\
\midrule
Media & Extiende el eje ``ensembles como inferencia'' de forma original (generaliza $\beta$ y $Z$), pero el contexto cuántico (cargas conservadas, integrabilidad) exige tiempo de exposición; útil como ejemplo avanzado si hay margen. \\
\bottomrule
\end{tabular}

\bigskip
\hrule
\bigskip

%=================================================================
\section{Mazzanti \& Romero (2020) --- Efficient Evaluation of the Partition Function of RBMs with AIS}

\subsection*{Referencia completa}
F. Mazzanti \& E. Romero (2020). \textit{Efficient Evaluation of the Partition Function of
RBMs with Annealed Importance Sampling}. arXiv:2007.11926 [cs.LG]. \cite{mazzanti2020}

\subsection*{Pregunta central}
¿Cómo elegir la distribución inicial y los parámetros de Annealed Importance Sampling (AIS)
para estimar eficientemente $\log Z$ de una RBM con pocas muestras y cadenas intermedias, sin
sacrificar precisión?

\subsection*{Aporte principal}
Muestra que partir de una distribución inicial $p_0(x)$ construida a partir de los sesgos
visibles del propio RBM (en vez de la distribución uniforme) reduce drásticamente el error de
la estimación de $\log Z$ vía AIS, tanto en modelos exactamente solubles como en problemas
reales (MNIST, Web-w6a), permitiendo monitorear $Z$ durante el entrenamiento con bajo costo
computacional.

\subsection*{Conexión con los ejes temáticos}
\begin{itemize}
  \item \textbf{Distribución de Boltzmann:} Indirecta --- el objeto estimado es la propia distribución de Boltzmann del RBM.
  \item \textbf{Pesos probabilísticos y EBM:} Sí --- usa los pesos y sesgos $w_{ij}, b_i$ del RBM para construir la distribución inicial óptima de AIS.
  \item \textbf{Función de partición $Z$:} Sí --- es el tema central y exclusivo del paper: estimación práctica y eficiente de $Z$.
  \item \textbf{Ensembles como inferencia:} No aplica directamente --- es un método numérico de muestreo, no un argumento de inferencia bayesiana.
\end{itemize}

\subsection*{Ecuaciones o resultados clave}
\begin{equation}
\log(Z_{\text{AIS}}) = \log\langle Z_n\rangle_s = \log\left[\frac{1}{N_s}\sum_i e^{s_i}\right] \approx \log(Z_n)
\end{equation}
\begin{equation}
\langle x \rangle_n = \frac{1}{e^{B}+1} \qquad \text{(condición de sesgo óptimo $B$ vía minimización KL)}
\end{equation}

\subsection*{Cita textual}
\begin{quote}
\textit{``a suitable starting probability distribution $p_0(x)$ can lead to a big improvement''}
\end{quote}

\subsection*{Relevancia para una presentación de 10 minutos}
\begin{tabular}{@{}p{2cm}p{11.5cm}@{}}
\toprule
\textbf{Nivel} & \textbf{Justificación} \\
\midrule
Alta & Aborda directamente el problema central del eje de la función de partición con una solución práctica y moderna; conecta teoría con implementación real. \\
\bottomrule
\end{tabular}

\bigskip
\hrule
\bigskip

%=================================================================
\section{Dawid \& LeCun (2023) --- Introduction to Latent Variable Energy-Based Models}

\subsection*{Referencia completa}
A. Dawid \& Y. LeCun (2023). \textit{Introduction to Latent Variable Energy-Based Models: A
Path Towards Autonomous Machine Intelligence}. Les Houches Summer School Lecture Notes;
Journal of Statistical Mechanics: Theory and Experiment. DOI:
\href{https://doi.org/10.1088/1742-5468/ad292b}{10.1088/1742-5468/ad292b}. arXiv:2306.02572.
\cite{dawid2023}

\subsection*{Pregunta central}
¿Cómo deben extenderse los EBM con variables latentes para soportar predicción bajo
incertidumbre y constituir una arquitectura de ``inteligencia artificial autónoma'' (JEPA) más
allá del aprendizaje supervisado o por refuerzo?

\subsection*{Aporte principal}
Introduce los EBM de variable latente, $F_w(x,y)=\min_z E_w(x,y,z)$ (o su versión
marginalizada vía energía libre), mostrando cómo minimizar/marginalizar sobre $z$ absorbe la
incertidumbre irreducible de la predicción. Propone la arquitectura JEPA jerárquica como camino
hacia sistemas de IA con ``sentido común''.

\subsection*{Conexión con los ejes temáticos}
\begin{itemize}
  \item \textbf{Distribución de Boltzmann:} Sí --- recupera explícitamente la distribución de Gibbs-Boltzmann como puente entre EBM y modelos probabilísticos.
  \item \textbf{Pesos probabilísticos y EBM:} Sí --- es una extensión moderna directa del marco EBM, añadiendo variables latentes.
  \item \textbf{Función de partición $Z$:} Sí --- discute la marginalización sobre $z$, formalmente idéntica a la energía libre $F=-kT\ln Z$.
  \item \textbf{Ensembles como inferencia:} Parcial --- la marginalización sobre variables latentes es análoga al ensemble canónico, aunque el paper no desarrolla explícitamente la lectura bayesiana.
\end{itemize}

\subsection*{Ecuaciones o resultados clave}
\begin{equation}
P(y\mid x) = \frac{e^{-\beta F(x,y)}}{\int dy'\, e^{-\beta F(x,y')}}
\end{equation}
\begin{equation}
F_w(x,y) = -\frac{1}{\beta}\log \int dz'\, e^{-\beta E_w(x,y,z')}
\end{equation}

\subsection*{Cita textual}
\begin{quote}
\textit{``the energy function is not the objective function to minimize''}
\end{quote}

\subsection*{Relevancia para una presentación de 10 minutos}
\begin{tabular}{@{}p{2cm}p{11.5cm}@{}}
\toprule
\textbf{Nivel} & \textbf{Justificación} \\
\midrule
Alta & Conecta los ejes de EBM y función de partición con la frontera actual de la investigación en IA (JEPA, LeCun), dando proyección y actualidad a la charla. \\
\bottomrule
\end{tabular}

\bigskip
\hrule
\bigskip

%=================================================================
\section{Carbone (2025) --- Hitchhiker's Guide on Energy-Based Models}

\subsection*{Referencia completa}
D. Carbone (2025). \textit{Hitchhiker's Guide on the Relation of Energy-Based Models with
Other Generative Models, Sampling and Statistical Physics: A Comprehensive Review}.
Transactions on Machine Learning Research, 05/2025. arXiv:2406.13661.
\href{https://openreview.net/forum?id=VTgixSbrJI}{OpenReview}. \cite{carbone2025}

\subsection*{Pregunta central}
¿Cómo se relacionan los EBM con otros paradigmas generativos (GANs, VAEs, normalizing flows), y
qué papel desempeñan los métodos de muestreo (MCMC) y la física estadística de no equilibrio en
su entrenamiento moderno?

\subsection*{Aporte principal}
Ofrece una revisión exhaustiva dirigida a físicos que traduce el lenguaje de los EBM al de la
física estadística, cubriendo desde la función de energía/partición hasta las técnicas de
entrenamiento ``sampling-free'' más recientes, y propone integrar herramientas de termodinámica
de no equilibrio (dinámica de Langevin, teoremas de fluctuación) en el entrenamiento de EBM.

\subsection*{Conexión con los ejes temáticos}
\begin{itemize}
  \item \textbf{Distribución de Boltzmann:} Sí --- es el punto de partida explícito de toda la revisión.
  \item \textbf{Pesos probabilísticos y EBM:} Sí --- eje central de la revisión completa.
  \item \textbf{Función de partición $Z$:} Sí --- discute extensamente su intratabilidad y los métodos para evitarla o estimarla (MCMC, sampling-free).
  \item \textbf{Ensembles como inferencia:} Parcial --- discute el lenguaje físico (Langevin, fluctuaciones) pero no enfatiza la lectura bayesiana del eje 4 tan explícitamente como Jaynes.
\end{itemize}

\subsection*{Ecuaciones o resultados clave}
\begin{equation}
\rho_\theta(x) = \frac{1}{Z_\theta} e^{-U_\theta(x)}, \qquad Z_\theta = \int e^{-U_\theta(x)}\,dx
\end{equation}
\begin{equation}
dX_t = -\nabla U(x)\,dt + \sqrt{2}\,dW_t \qquad \text{(dinámica de Langevin para muestreo)}
\end{equation}

\subsection*{Cita textual}
\begin{quote}
\textit{``highlighting the significance of energy functions and partition functions''}
\end{quote}

\subsection*{Relevancia para una presentación de 10 minutos}
\begin{tabular}{@{}p{2cm}p{11.5cm}@{}}
\toprule
\textbf{Nivel} & \textbf{Justificación} \\
\midrule
Media & Es una síntesis moderna excelente que vincula los cuatro ejes con el estado del arte, pero su amplitud (GANs, flows, no equilibrio) excede lo abordable en 10 minutos; mejor como mapa de referencia que como contenido a desarrollar. \\
\bottomrule
\end{tabular}

\bigskip
\hrule
\bigskip

%=================================================================
\section{Hohm (2026) --- Lecture Notes on Statistical Physics and Neural Networks}

\subsection*{Referencia completa}
O. Hohm (2026). \textit{Lecture Notes on Statistical Physics and Neural Networks}.
arXiv:2605.06394 [cond-mat.dis-nn]. \cite{hohm2026}

\subsection*{Pregunta central}
¿Cómo puede presentarse la física estadística clásica (distribución de Boltzmann-Gibbs, modelo
de Ising, grupo de renormalización) como base pedagógica directa para entender las máquinas de
Boltzmann y el aprendizaje profundo moderno?

\subsection*{Aporte principal}
Construye un puente didáctico explícito entre el grupo de renormalización (RG) y el
entrenamiento de RBM, mostrando que ``integrar'' las neuronas ocultas de una RBM es
formalmente el mismo procedimiento que integrar grados de libertad en RG. Cierra con una
introducción a modelos de lenguaje de gran escala desde la óptica de la física estadística.

\subsection*{Conexión con los ejes temáticos}
\begin{itemize}
  \item \textbf{Distribución de Boltzmann:} Sí --- eje central del capítulo 2: deriva la distribución de Boltzmann-Gibbs para espines de Ising.
  \item \textbf{Pesos probabilísticos y EBM:} Sí --- define la energía de la RBM y su regla de aprendizaje en términos de $w_{i\alpha}$.
  \item \textbf{Función de partición $Z$:} Sí --- muestra cómo $Z$ se cancela al calcular probabilidades condicionales $P(h\mid v)$, conectando con el ``integrar variables'' del RG.
  \item \textbf{Ensembles como inferencia:} Parcial --- no usa explícitamente el lenguaje bayesiano, pero la analogía RG$\leftrightarrow$aprendizaje profundo es conceptualmente afín al eje 4.
\end{itemize}

\subsection*{Ecuaciones o resultados clave}
\begin{equation}
E(v,h) = -\sum_{i,\alpha} w_{i\alpha}\, v_i h_\alpha - \sum_i b_i v_i - \sum_\alpha c_\alpha h_\alpha
\end{equation}
\begin{equation}
\Delta w_{i\alpha} = \varepsilon\left(\langle v_i h_\alpha\rangle_{\text{data}} - \langle v_i h_\alpha\rangle_{B}\right)
\end{equation}

\subsection*{Cita textual}
\begin{quote}
\textit{``hidden neurons are integrated out as in the renormalization group''}
\end{quote}

\subsection*{Relevancia para una presentación de 10 minutos}
\begin{tabular}{@{}p{2cm}p{11.5cm}@{}}
\toprule
\textbf{Nivel} & \textbf{Justificación} \\
\midrule
Alta & Material pedagógico reciente que conecta explícitamente los cuatro ejes con el grupo de renormalización y el deep learning; ancla la charla en un marco didáctico único. \\
\bottomrule
\end{tabular}

\bigskip
\hrule
\bigskip

%=================================================================
\section{Merhav (2026) --- Grouped Reverse Importance Sampling for the Partition Function}

\subsection*{Referencia completa}
N. Merhav (2026). \textit{Grouped Reverse Importance Sampling for the Partition Function}.
arXiv:2606.26748 [cs.IT]. \cite{merhav2026}

\subsection*{Pregunta central}
¿Puede reducirse el error cuadrático medio (MSE) en la estimación de la función de partición
$Z(\beta)$ agrupando muestras de la distribución de Boltzmann y aplicándoles un peso conjunto,
en vez de ponderar cada muestra por separado (RIS estándar)?

\subsection*{Aporte principal}
Introduce el muestreo de importancia inversa agrupado (GRIS) en tres variantes de creciente
sofisticación (NOL, FSW, VSW), demuestra que el peso conjunto óptimo depende solo de la energía
total del grupo, y muestra reducciones del MSE de 20--65\% (NOL) y de hasta un 30--40\%
adicional (VSW) frente al RIS estándar ($k=1$).

\subsection*{Conexión con los ejes temáticos}
\begin{itemize}
  \item \textbf{Distribución de Boltzmann:} Sí --- el método estima $Z(\beta)$ de $p(x)=e^{-\beta U(x)}/Z(\beta)$, la distribución de Boltzmann continua.
  \item \textbf{Pesos probabilísticos y EBM:} No aplica directamente --- es un método estadístico de muestreo, no involucra pesos aprendibles de una red.
  \item \textbf{Función de partición $Z$:} Sí --- es el objeto exclusivo de estudio del paper.
  \item \textbf{Ensembles como inferencia:} Indirecta --- vincula el problema con la evaluación de evidencia bayesiana y cantidades termodinámicas, sin desarrollar el diccionario completo del eje 4.
\end{itemize}

\subsection*{Ecuaciones o resultados clave}
\begin{equation}
Z(\beta) = \int_{\mathcal X} e^{-\beta U(x)}\,dx, \qquad p(x) = \frac{e^{-\beta U(x)}}{Z(\beta)}
\end{equation}
\begin{equation}
\ln \hat Z(\beta) = \ln M - \ln\left[\frac{1}{n}\sum_i m(x_i)\,e^{\beta U(x_i)}\right]
\end{equation}

\subsection*{Cita textual}
\begin{quote}
\textit{``a joint weight function allows a significantly richer freedom for optimization''}
\end{quote}

\subsection*{Relevancia para una presentación de 10 minutos}
\begin{tabular}{@{}p{2cm}p{11.5cm}@{}}
\toprule
\textbf{Nivel} & \textbf{Justificación} \\
\midrule
Baja & Es un desarrollo metodológico reciente y muy técnico (estadística de muestreo) centrado exclusivamente en el eje 3; demasiado especializado para 10 minutos, aunque sirve como nota de ``investigación activa'' al cierre. \\
\bottomrule
\end{tabular}

\bigskip
\hrule
\bigskip

%=================================================================
\section*{Referencias}

\begin{thebibliography}{99}

\bibitem{jaynes1957a}
Jaynes, E. T. (1957). Information Theory and Statistical Mechanics. \textit{Physical Review},
106(4), 620--630. DOI: 10.1103/PhysRev.106.620.

\bibitem{jaynes1957b}
Jaynes, E. T. (1957). Information Theory and Statistical Mechanics II. \textit{Physical
Review}, 108(2), 171--190. DOI: 10.1103/PhysRev.108.171.

\bibitem{ackley1985}
Ackley, D. H., Hinton, G. E. \& Sejnowski, T. J. (1985). A Learning Algorithm for Boltzmann
Machines. \textit{Cognitive Science}, 9(1), 147--169. DOI: 10.1207/s15516709cog0901\_7.

\bibitem{lecun2006}
LeCun, Y., Chopra, S., Hadsell, R., Ranzato, M. \& Huang, F. J. (2006). A Tutorial on
Energy-Based Learning. En: \textit{Predicting Structured Data}. MIT Press.

\bibitem{cheng2018}
Cheng, S., Chen, J. \& Wang, L. (2018). Information Perspective to Probabilistic Modeling:
Boltzmann Machines versus Born Machines. \textit{Entropy}, 20(8), 583. DOI: 10.3390/e20080583.

\bibitem{puskarov2018}
Pu\v{s}karov, T. \& Cort\'es Cubero, A. (2018). Machine Learning Algorithms Based on
Generalized Gibbs Ensembles. arXiv:1804.03546 [cond-mat.stat-mech].

\bibitem{mazzanti2020}
Mazzanti, F. \& Romero, E. (2020). Efficient Evaluation of the Partition Function of RBMs with
Annealed Importance Sampling. arXiv:2007.11926 [cs.LG].

\bibitem{dawid2023}
Dawid, A. \& LeCun, Y. (2023). Introduction to Latent Variable Energy-Based Models: A Path
Towards Autonomous Machine Intelligence. \textit{Journal of Statistical Mechanics: Theory and
Experiment}. DOI: 10.1088/1742-5468/ad292b.

\bibitem{carbone2025}
Carbone, D. (2025). Hitchhiker's Guide on the Relation of Energy-Based Models with Other
Generative Models, Sampling and Statistical Physics: A Comprehensive Review.
\textit{Transactions on Machine Learning Research}, 05/2025. arXiv:2406.13661.

\bibitem{hohm2026}
Hohm, O. (2026). Lecture Notes on Statistical Physics and Neural Networks. arXiv:2605.06394
[cond-mat.dis-nn].

\bibitem{merhav2026}
Merhav, N. (2026). Grouped Reverse Importance Sampling for the Partition Function.
arXiv:2606.26748 [cs.IT].

\end{thebibliography}

\end{document}
